% \textcolor{red}{JUST ADDING A SKETCH. FILL THIS OUT. MAKE THIS FLOW BETTER, MORE COHERENT, AND ADD CITATIONS WHEREVER NEEDED.}

% Talk about LLMs being used in real world and these include settings where safety is critical. This necessitates LLMs that produce helpful responses that are safe (i.e no toxicity, misinformation)
Large Language Models (LLMs) are now used across a wide range of applications, so it is important that their outputs are helpful, reliable, and safe. This need is especially acute in high-stakes domains such as legal reasoning (\cite{katz2024gpt}), medical consultation (\cite{yang2022large,moor2023foundation}), and educational support (\cite{kasneci2023chatgpt,kung2023performance}), where harmful generations—including toxicity and misinformation—can have serious consequences (\cite{gehman2020realtoxicityprompts,weidinger2021ethical,ganguli2022red}).

Broadly, we want LLMs to be \textit{helpful} and aligned with human preferences, while being \textit{harmless}, i.e. prevent toxic outputs. But there may be scenarios when users may request assistance with potentially harmful activities (\cite{Bai2022ConstitutionalAH,glaese2022improving}). Thus, \emph{harmlessness} and \emph{helpfulness} can conflict. The standard technique for aligning LLMs with human preferences, Reinforcement Learning from Human Feedback (RLHF), typically optimizes a single reward model for both objectives (\cite{Ouyang2022TrainingLM,bai2022traininghelpfulharmlessassistant}), or heuristically mixes separate reward models (\cite{glaese2022improving,touvron2023llama,mu2024rule}). This creates trade-offs: stronger harmlessness can lead to refusals, while stronger helpfulness can increase unsafe outputs (\cite{bai2022traininghelpfulharmlessassistant}). Recent work addresses this by decoupling preference data and enforcing harmlessness as a constraint, an approach known as Safe RLHF (\cite{dai2023safe}). However, Safe RLHF typically typically constrains the expected cost of a policy, providing no guarantees on worst-case or tail outcomes.  This is inadequate in high-stakes deployments where tail risk is critical — such as code generation \cite{pearce2021asleepkeyboardassessingsecurity} and open-ended dialogue where rare but severe harms including toxic generations and private data leakage can impact users
\citep{perez2022redteaminglanguagemodels, gehman2020realtoxicityprompts}.

%To overcome this limitation, Safe RLHF modifies both the reward modeling and reinforcement learning phases, explicitly introducing a safety constraint to minimize harmfulness while preserving helpfulness.

% Helpfullness and Safety can conflict with each other. Talk about seperating these objectives, framing the problem as a constrained optimization problem, which was used in Safe-RLHF.

%However, , yet standard
% FSD is a principled approach to compare distributions and better handles uncertainity. FSD is partial ordering, hence making it harder to optimize

\yc{Also, I changes the flow of the old paragraph, and incorporated scott's comments. Lmk if this seems good or not}\raj{we already talk about expected cost contraints above. Moved the first line to the pargarph above since it is more expressive.}
%Safe alignment methods typically constrain the \textit{expected} cost of a policy, providing no guarantees on worst-case or tail outcomes.
Thus, we argue that a stronger notion of constraining the cost of the policy is needed: the cost distribution of the learned policy should be \emph{stochastically smaller} than that of the reference policy, not merely cheaper on average. \raj{added a linbe which explains what it means to be stcoshstically smaller} This means that the learned policy should assign less probability to high-cost outcomes across the distribution, not just reduce the mean cost. This motivates \textbf{Risk-sensitive Alignment via Dominance (RAD)}, which enforces a First-Order Stochastic Dominance (FSD) constraint on the cost distribution of the learned policy relative to a reference policy.

A key insight of RAD is that by reweighting quantiles of the FSD objective, we recover Spectral Risk Measures (SRMs) — a family of risk measures defined as a weighted integral of the cost quantiles,
$\int w(q) Q_X(q) dq$, where $Q_X(q)$ is the quantile function of the cost distribution and the weighting function $w(q)$ allows
practitioners to express diverse risk preferences over the cost distribution. For instance, concentrating weight on upper quantiles recovers tail-sensitive measures such as CVaR, while uniform weighting recovers the mean — both expressible within our unified framework. This is practically valuable: a medical deployment may demand near-zero tolerance for harmful outputs, while a general assistant
may tolerate a more permissive tradeoff, and both are accommodated by choosing $w$ appropriately.

Optimizing FSD constraints directly is challenging, so we relax them to an asymmetric FSD-violation surrogate that aggregates positive quantile gaps (see \Cref{eq:fsd_loss}). We interpret stochastic
dominance through Optimal Transport (OT) and derive an efficient REINFORCE-style policy gradient using Sinkhorn iterations, yielding an end-to-end differentiable approach. Empirically, RAD yields models that are more robust to safety violations while achieving competitive helpfulness across multiple spectral risk measures relative to Safe RLHF.
% and under several spectral risk measures achieve simultaneous improvements in both safety and helpfulness. 
\yc{IS this ok, given the results?}\raj{modified a bit}

% In this work, we propose \textbf{Risk-sensitive Alignment via Dominance (RAD)}, which performs safe alignment by enforcing a First-Order Stochastic Dominance (FSD) constraint that makes the learned policy's costs stochastically smaller than a reference policy's costs.
% Optimizing FSD constraints is challenging, so we relax them to an asymmetric FSD-violation surrogate that aggregates positive quantile gaps (see \Cref{eq:fsd_loss}).
% We interpret stochastic dominance through Optimal Transport (OT) and propose an efficient REINFORCE-style algorithm. Using Sinkhorn iterations, we derive a policy-gradient expression for the FSD objective, yielding an end-to-end differentiable approach to optimize FSD-based constraints or objectives. We further show that reweighting quantiles recovers Spectral Risk Measures (SRMs): imposing FSD constraints with a given weighting yields, at convergence, a reduction in the corresponding SRM whenever reverse violations vanish.
% % imposing FSD loss constraints with a given weighting guarantees, at convergence, a reduction in the corresponding SRM relative to the reference policy's cost distribution.
% This lets practitioners tune the risk profile by choosing the weighting function, providing a unified approach to safe alignment. Experimentally, SARD yields models that are more robust to safety violations while achieving competitive helpfulness compared to expected-constraint methods such as Safe-RLHF.

% \aj{Need to add a para about how the paper is structured and one or two lines about the experiments and observations.}

\aj{verify}
Our main contributions are: (i) we formulate safe alignment as a dominance-constrained objective over the full cost distribution rather than only its expectation, (ii) we derive a RAD policy-gradient estimator using entropic OT and Sinkhorn iterations, making optimization of the objective end-to-end differentiable, (iii) we connect quantile reweighting in our framework to SRMs, enabling controllable risk-sensitive alignment profiles, and (iv) through extensive experiments, we show that RAD improves robustness to safety violations while maintaining competitive helpfulness relative to expected cost based baselines. \aj{extra} The rest of the paper is organized as follows: \Cref{sec: background} reviews Safe RLHF, stochastic dominance, and SRMs; \Cref{sec: method} presents RAD and its optimization procedure; \Cref{sec: results} reports empirical evaluations safety--helpfulness trade-offs; and \Cref{sec: relatedwork,sec: conclusion} is related work and conclusion respectively.
\aj{Can write one or two lines highlighting the conclusion of experiments.}\raj{changed second point. Replacing unregularized with regaularized OT is not really an original contribution and pretty standard. But we show how to estimate the gradients which could be a new contribution} \yc{Add experimental conclusions as a point in our contribution} \raj{added experimental contributional.}
